\documentclass[11pt,a4paper]{article}

\usepackage[utf8]{inputenc}
\usepackage[T1]{fontenc}
\usepackage{lmodern}
\usepackage{amsmath,amssymb}
\usepackage{graphicx}
\usepackage{booktabs}
\usepackage{siunitx}
\usepackage{microtype}
\usepackage[margin=1in]{geometry}
\usepackage{hyperref}
\usepackage[capitalise,noabbrev]{cleveref}

\title{ Stabilizing Cortical Wetware by Natural Algorithm Optimization }
\author{Eliza Kishan\\
	\small BSc '27, Department of Computer Science and Department of Biology, McGill University \\
	\small \texttt{eliza.kishan@mail.mcgill.ca}}
\date{\today}

\begin{document}

\maketitle

\begin{abstract}
Briefly summarize the research question, methods, principal results, and conclusions.
\end{abstract}

\section{Introduction}
Cortical wetware is the use of human neurons, typically derived from iPSC stem cells, due to recent developments in both ease and efficiency in the chemical
method \cite{ipscneuron}.

\section{Related Work}


\section{Methods}
\subsection{Data, preprocessing and in-silico model definition}

The entire analysis is performed using a 1000-neuron connected cell-network
within an extracellular matrix, all simulated using NetPyNE and NEURON \cite{neuron} \cite{netpyne}. The set-up for the cell
was based upon the need to develop close-to-reality exchange of ions and GABA receptor dynamics which induce the unstable activity
we wish to resolve. 

Seizure models were deried from single-cell models, which are then applied to the wider system. In order to systematize the single-cell models,
NetPyNE was parameterized to adopt <fill in here>.

\subsection{Analysis}
The "accuracy" or performance of the model was evaluated based on specific parameters.


\subsection{Algorithms}
\subsubsection{Neural Nets}
\subsubsection{Ant Colony Optimization}
\subsubsection{Bee Colony Optimization}
\subsubsection{Artificial Immune Systems}
\subsubsection{Firefly Algorithms}
\subsubsection{Particle Swarm Optimization \& Genetic Algorithms}


\section{Results}
Present the findings without interpretation. Use tables and figures as needed (e.g., \cref{fig:example}).

\begin{table}[ht]
	\centering
	\caption{Example results table.}
	\label{tab:results}
	\begin{tabular}{lcc}
		\toprule
		Method & Metric 1 & Metric 2 \\
		\midrule
		Baseline & -- & -- \\
		Proposed & -- & -- \\
		\bottomrule
	\end{tabular}
\end{table}

\begin{figure}[ht]
	\centering
	% \includegraphics[width=0.8\textwidth]{figures/example}
	\caption{Example figure caption.}
	\label{fig:example}
\end{figure}

\section{Discussion}
Interpret the results, compare them with prior work, and discuss implications.

\section{Limitations}
Describe assumptions, sources of uncertainty, and limitations of the study.

\section{Conclusion}
Summarize the main findings and identify promising directions for future work.

\section*{Acknowledgments}
This work was conducted as a final project for COMP 599, a topics course at the Dept. of Computer Science, McGill University. I thank Dr. Waldispuhl for
his support and guidance during this project. 

\subsection*{Code Availability}
Code of the wetware simulator and of the different stabilization algorithms is available publicly at the GitHub link: https://github.com/StarlightAbove/wetware-stability. All work and code is protected by
the MIT License.

\subsection*{AI Declaration}
Some of the codebase was developed and debugged via the use of Claude Opus 5.5 at high effort, a LLM from Anthropic PBC, San Francisco. Some literature search was conducted by 
Claude Sonnet 5.5 at high effort. All writing and conception was done by hand.

\bibliographystyle{plain}
\begin{thebibliography}{9}

\bibitem{ipscneuron}
Zhang, Yingsha et al. “Rapid single-step induction of functional neurons from human pluripotent stem cells.” Neuron vol. 78,5 (2013): 785-98. doi:10.1016/j.neuron.2013.05.029.

\bibitem{neuron}
Hines, Michael, Thomas Carnevale, and Robert A. McDougal. "NEURON Simulation Environment." Encyclopedia of Computational Neuroscience, edited by Dieter Jaeger and Ranu Jung, Springer, 2022, doi: 10.1007/978-1-0716-1006-0\_795.

\bibitem{netpyne}
Dura-Bernal, Salvador et al. “NetPyNE, a tool for data-driven multiscale modeling of brain circuits.” eLife vol. 8 e44494. 26 Apr. 2019, doi:10.7554/eLife.44494.

\end{thebibliography}

\end{document}
